\uuid{Vatm}
\exo7id{5101}
\titre{exo7 5101}
\auteur{rouget}
\organisation{exo7}
\datecreate{2010-06-30}
\isIndication{false}
\isCorrection{true}
\chapitre{Continuité, limite et étude de fonctions réelles}
\sousChapitre{Limite de fonctions}
\cpge{mpsi-mp2i.fonctions-valeurs-reelles-complexes.c, pcsi.techniques-calcul-fonctions.c, ptsi.techniques-calcul-fonctions.c, mpsi-mp2i.limites-continuite.a, pcsi.limites-continuite.a, ptsi.limites-continuite.a, bcpst1.fonctions-reelles-usuelles.c, bcpst1.limites-continuite.a}
\module{Analyse}
\niveau{L1}
\difficulte{}

\contenu{
\question{
Trouver $\lim_{x\rightarrow +\infty}\frac{(x^x)^x}{x^{(x^x)}}$.
}
\reponse{
Pour $x>0$, $(x^x)^x=e^{x\ln(x^x)}=e^{x^2\ln x}$ et $x^{(x^x)}=e^{x^x\ln x}$. Par suite,

$$\forall x>0,\;\frac{(x^x)^x}{x^{(x^x)}}=\text{exp}(\ln x(x^2-x^x)).$$
Or, $x^2-x^x=-x^x(1-x^{2-x})=-e^{x\ln x}(1-e^{(2-x)\ln x})$. Quand $x$ tend vers $+\infty$, $(2-x)\ln x$ tend vers
$-\infty$. Donc, $1-e^{(2-x)\ln x}$ tend vers $1$ puis $x^2-x^x$ tend vers $-\infty$. Mais alors, $\ln x(x^2-x^x)$ tend
vers $-\infty$, puis $\frac{(x^x)^x}{x^{(x^x)}}=\mbox{exp}(\ln x(x^2-x^x))$ tend vers $0$.

\begin{center}
\fbox{
$\lim_{x\rightarrow +\infty}\frac{(x^x)^x}{x^{(x^x)}}=0.$
}
\end{center}
}
}
